\scanpage{60}

\textcircled{f}
\[
\begin{aligned}
X,Y\in\unsure{\mathcal E|\mathscr E}_{\chi_!}&\Rightarrow X\times Y\in\mathcal E_{\chi_!} &&\text{et on a} &\chi_!(X\times Y)&=\chi_!(X)\chi_!(Y)\\
X,Y\in\mathcal E_\chi&\Rightarrow X\times Y\in\mathcal E_\chi &&\text{et on a} &\chi(X\times Y)&=\chi(X)\chi(Y).
\end{aligned}
\]

(La première assertion résulte de Künneth pour $H^*_!(X\times Y,k)$, la deuxième de Künneth pour $H_*(X\times Y,k)$, dont le dual est $H^*(X\times Y,k)$, donc qui permet de tester l'appartenance à $\mathcal E_\chi$ et la valeur de $\chi$ $\ldots$)

\textcircled{g} Soit $f:X\unsure{-|\longrightarrow}Y$ fibration loc. triviale de fibre $F$.

Si $Y,F\in\mathcal E_{\chi_!}$ ; alors $X\in\mathcal E_{\chi_!}$ et on a
\[
\chi_!(X)=\chi_!(Y)\chi_!(F).
\]
\margin{avec léger grain de sel}% slanted

Énoncé analogue pour \inserted{$Y,F\in$} $\mathcal E_\chi$ et \deleted{\unsure{$\chi(X)=$}}
\[
\chi(X)=\chi(Y)\chi(F)
\]

\uline{Dém.} On a
\[
H^*_!(X)\Leftarrow E_2^{pq}=H^p_!\bigl(Y,\underbrace{R^q f_!(k_X)}_{\substack{\text{syst. local sur }Y\\\text{à fibres }\sim H^q_!(F)}}\bigr)
\]
donc \uline{sous réserve} qu'on sache que $E_2^{**}$ de dim finie, on trouve $\deleted{\unsure{\chi}}X\in\mathcal E_\chi$ et
\[
\begin{aligned}
\chi(X)&=\sum_{pq}(-1)^{p+q}\operatorname{rg}_k H^p_!(Y,R^q f_!(k_X))\\
&=\sum_q(-1)^q\sum_p\underbrace{(-1)^p\operatorname{rg}_k H^p_!(Y,R^q f_!(k_X))}_{b^q_!(F)\chi_!(Y)}\\
&=\chi_!(Y)\deleted{\unsure{\chi}}\sum_q(\deleted{\unsure{+}}-1)^q b^q_!(F)=\chi_!(Y)\chi_!(F)
\end{aligned}
\]

Il faut prouver

\uline{Lemme} Soit $\mathcal R$ syst. local de $k$-vectoriels sur \deleted{\unsure{$Y$|$X$} \unsure{loc. cpt}}, $Y\in\mathcal E_{\chi_!}$, à fibres de rang \inserted{fini} $b$. Alors
\[
\chi_!(Y,\mathcal R)\overset{\mathrm{déf}}{=}\sum_p(-1)^p H^p(Y,\mathcal R)
\]
existe (i.e. $\unsure{\sum|\bigoplus}_p H^p(Y,\mathcal R)$ de dim finie sur $k$) et on a
\[
\chi_!(Y,\mathcal R)=b\chi_!(Y)
\]

Marche tout au moins si $Y$ \deleted{triangulable} \inserted{admet filtration finie} avec $Y_i-Y_{i-1}$ somme finie de cellules \inserted{$Z_{ij}$} $\simeq R^{\unsure{d|\delta}_{ij}}$, car
\margin{\textcircled{\unsure{gr.|gn|gr}}}
On aura
\[
\begin{aligned}
\chi_!(Y,\mathcal R)&=\sum_{ij}\chi_!(Z_{ij},\underset{(\text{ct. sur }Z_{ij})}{\mathcal R})=\sum_{ij}b\chi_!(Z_{ij})\\
&=b\sum\chi_!(Z_{ij})\unsure{\mathpunct{:}|=}b\chi_!(Y).
\end{aligned}
\]

% archive sheet number: 59

